\section{Preliminaries}
\label{sec:prelims}
\begin{figure}[H]
\small
\begin{tabularx}{\textwidth}{rll}
   $n, m \in$
   &
   $\Int$
   &
   integer
   \\
   $i, j \in$
   &
   $\set{n \in \Int \mid n \numgeq 1}$
   &
   positive integer
   \\
   $z \in$
   &
   $\Str$
   &
   finite sequence of characters
   \\
   $x, y, f, \oplus \in$
   &
   $\Var$
   &
   variable
   \\
   $c \in$
   &
   $\Ctr$
   &
   constructor name
   \\
   $D \in$
   &
   $\DataType$
   &
   datatype name
\end{tabularx}
\caption{Basic elements}
\label{fig:basic-elements}
\end{figure}


Variables include a distinguished variable $\varAnon$ called the \emph{anonymous}
variable.

For any constructor $c$ we write $\datatype{\upSigma}{c}$ to mean either the datatype that $c$ is a member of
or for the type signature of $c$ as a function. The choice of meaning should be obvious from context.

\subsection{Foreign functions}

The language is parameterised by a finite map $\Phi: \Var \to \Type^+ \times \Type$ from variables to foreign
function signatures $(\seq{A}, B)$. By convention $f$ and $\oplus$ ranges over $\dom{\Phi}$. Foreign functions
\emph{per se} are not first-class and calls $\exForeignApp{f}{\seq{e}}$ are saturated, i.e.~$\length{\seq{e}}
= \length{\seq{A}}$. First-class foreign functions are provided via a top-level environment which maps every
$\bind{f}{(\seq{A}, B)} \in \Phi$ to the closure
$\exClosure{\seqEmpty}{\seqEmpty}{\elimVar{x_1}{\elimVar{\iter}{\elimVar{x_i}{\exForeignApp{f}{\seq{x}}}}}}$
where $\length{\seq{x}} = \length{\seq{A}}$. Although saturated calls are technically part of the surface and
core languages, user programs by convention use the first-class representation.

\subsection{Continuations and Eliminators}

We explicitly represent the notion of a continuation in the language, represented by the metavariable $\upkappa$.
A continuation can be one of either a term or an \textit{eliminator}. Eliminators are a trie-like structure which
in the core language allow us to use a uniform syntax and semantics for pattern-matching. An eliminator $\upsigma$
is said to be \textit{partial} if we do not require that every possible case of pattern matching appears within it.

Eliminators are polymorphic in the type of the object they return, which we call the continuation type. An eliminator
can either return an eliminator, or an expression we usually refer to as the body.
We overload the use of $\elimTyArrow$, as it appears in both the syntax of eliminators, and the typing judgement for
continuation types. The meaning should be clear from context.
\begin{figure}
   \begin{minipage}[t]{0.48\textwidth}
      \begin{tabularx}{\textwidth}{rL{2.7cm}L{3cm}}
          &\textit{Eliminator}&
          \lowlight{$\Elim\;\upGamma\,A\;K$}
          \\
          $\upsigma, \uptau ::=$
          &
          $\elimVar{x}{\upkappa}$
          &
          variable
          \\
          &
          $\elimRec{\seq{x}}{\upkappa}$
          &
          record
          \\
          &
          $\elimConstr{\seq{\elimBind{c}{\upsigma}}}$
          &
          data value
          \\[2mm]
          &\textit{Recursive functions}&
          \lowlight{$\RecDefs\;\upGamma\;\upGamma'$}
          \\
          $\uprho ::=$
          &
          $\set{\seq{\bind{x}{\upsigma}}}$
          &
          \\[2mm]
          &\textit{Value}&
          \lowlight{$\Val{A}$}
          \\
          $u, v ::=$
          &
          $\exInt{n}$
          &
          integer
          \\
          &
          $\exStr{z}$
          &
          string
          \\
          &
          $\exRec{\seq{\bind{x}{v}}}$
          &
          record
          \\
          &
          $\exDict{\seq{\bind{s}{v}}}$
          &
          dictionary
          \\
          &
          $\exConstr{c}{\seq{v}}$
          &
          data value
          \\
          &
          $\exClosure{\upgamma}{\uprho}{\upsigma}$
          &
          closure
          \\[0mm]
      \end{tabularx}
   \end{minipage}
   \caption{Syntax of Eliminators}
\end{figure}

\subsection{Important Relations}
Later on in the paper, we make use of a large number of definitions in the form of relations. These relations 
we have defined are used for pretty much all things we use in the draft. We have very standard relations
for evaluation ($\evalFwdS,\,\evalSugS$), desugaring ($\desugarFwdS$) and pattern matching ($\matchFwdS$).

We also have some more esoteric relations, for joining eliminators ($\disjjoin$), and for closing environments
($\closeDefsS$). We give a short account of each relation here.

% \subsubsection{$\evalFwdS$}
% \subsubsection{$\evalSugS$}
% \subsubsection{$\desugarFwdS$}
% \subsubsection{$\protect\closeDefsS$}
% \subsubsection{$\matchFwdS$}
% \subsubsection{$\disjjoin$}
\clearpage
